\documentclass[10pt,a4paper]{amsart}
\usepackage[margin=19mm]{geometry}
\usepackage{amsmath,amssymb,mathtools}
\usepackage[scr=rsfs,cal=euler]{mathalfa}
\usepackage{xfrac}
\usepackage[unicode,hidelinks,bookmarks=false]{hyperref}
\newcommand{\ie}{i.e.\ }
\newcommand{\cf}{cf.\ }
\newcommand{\resp}{respectively\ }
\newcommand{\Subsection}{\subsection}
\providecommand{\nobreakdash}{\nobreak\hbox{-}}
\setlength{\emergencystretch}{2em}
\multlinegap0pt
\usepackage{amssymb}
\usepackage{mathtools}
\usepackage{xcolor}
\usepackage{dsfont}
\newtheorem{lemma}{Lemma}
\def\N{\mathbb{N}}
\def\om{\omega}
\title{Effective quenched linear response for random dynamical systems}
\author{Davor Dragicevic, Yeor Hafouta}
\date{Source: arXiv:2403.04907v3 (CC BY 4.0)}
\begin{document}
\maketitle

\noindent\emph{Source and licence.} Effective quenched linear response for random dynamical systems. Version arXiv:2403.04907v3 (CC BY 4.0), distributed by arXiv under the Creative Commons Attribution 4.0 International licence, http://creativecommons.org/licenses/by/4.0/, as recorded in the arXiv licence metadata for this exact version and shown on the abstract page https://arxiv.org/abs/2403.04907v3. Full licence notice: \texttt{licenses/arxiv-2403.04907-CC-BY-4.0.txt}.\par\smallskip

\noindent\textit{Editorial scope.} Counted unit: the article's lemma 4 (source0.tex lines 1747--1772). The article's complete original proof of it is delivered with it (source0.tex lines 2457--2670, one \texttt{proof} environment of 214 lines, which the article places away from the statement). No mathematics is written, repaired or re-derived: the statement and the proof are the article's own lines, in the article's own notation, and no retained line is substituted or re-flowed.  Retained uncounted as the source-local context the proof needs: lines 266--272; lines 1433--1435; lines 1717--1719.  Nothing was omitted from any retained span, so the package carries no omission record.  No retained line contains a citation, so the wrapper carries no bibliography and no bibliography entry.  No retained line contains a figure environment, an image-inclusion command or drawing code, so no image is delivered and the package declares no asset dependency.  Editorial formatting only, in addition: an AMS \texttt{amsart} preamble that restates the article's own declarations and macros used by the retained lines, namely the environments \texttt{lemma} on separate counters, together with the standard packages those lines need.  The retained environments print in this document as Lemma 1 in order of appearance, whereas the article prints the same items with the numbering of its own full document.  The complete native source of the article as distributed in the arXiv e-print for this exact version is bundled unchanged as \texttt{sources/155745/source0.tex}.
\begin{lemma}\label{B lemma}
Suppose that  $B\in L^q(\Omega,\mathcal F,\mathbb P)$ for some $q>0$.  Then, for every sequence of positive numbers $(a_n)_{n\in \N}$ such that $\sum_{n\geq 1}a_n^q<+\infty$, there is a random variable $R\in L^q(\Omega,\mathcal F,\mathbb P)$ such that
\begin{equation}\label{BR}
B(\sigma^{-n}\omega)\leq R(\omega)a_n^{-1},
\end{equation}
for $\mathbb P$-a.e. $\omega \in \Omega$ and every $n\in \N$. In particular, for every $\delta>0$ there exists $R_\delta\in L^q(\Omega,\mathcal F,\mathbb P)$ such that $B(\sigma^{-n}\om)\leq R_\delta(\om)n^{1/q+\delta}$ for $\mathbb P$-a.e. $\om \in \Omega$ and $n\in \N$.
\end{lemma}

\begin{equation}\label{gam dec}
\max\left(\prod_{j=0}^{n-1}\gamma_{\sigma^j\omega}^{-1},\prod_{j=1}^{n-1}\gamma_{\sigma^{-j}\omega}^{-1}\right)\leq  E_\omega n^{-a_0}.    
\end{equation}

\begin{equation}\label{Pre Adler}
\max(\|D^2 y_{\omega,\varepsilon}\|_\infty, \|D^3 y_{\omega,\varepsilon}\|_\infty, \|D^4 y_{\omega,\varepsilon}\|_\infty)\leq c(\omega).
\end{equation}

\begin{lemma}\label{Lemma 4}
Suppose that $\gamma_\omega\geq1$ and that $\gamma_\omega\in L^q$ for some $q>1$. Let   \eqref{Pre Adler} hold and
assume also that $a_0>1+\frac{1}{p}+\frac1{q}$, where $a_0$ comes from \eqref{gam dec} and $p$ comes from \eqref{Pre Adler}.
Then, for every $\delta>0$ (small enough) there exists a random variable $C_\omega\geq 1$ such that
for all $n$, $\varepsilon\in I$ and every inverse branch $y$ of $T_{\sigma^{-n}\omega,\varepsilon}^n$ we have 
\begin{equation}\label{11}
\|D^2y\|_\infty\leq C_\omega n^{-\eta}    
\end{equation}
where $\eta=a_0-1-1/s-\delta$,
$s$ is given by $1/s=1/p+1/q$,  $\omega\mapsto C_\omega\in L^{t}(\Omega,\mathcal F,\mathbb P)$, and $t$ is given by $\frac{1}{t}=\frac{1}{q_0}+\frac{1}{p}+\frac{1}{q}$ ($a_0$ and $q_0$ come from \eqref{gam dec}). 

Moreover, if $a_0>2+1/s$ then
\begin{equation}\label{12}
\|D^3y\|_\infty\leq A_\om n^{-\zeta}    
\end{equation}
where $\zeta=a_0-2-1/s-\delta$, 
and $A_\om\geq 1$ is a random variable such that $\om\mapsto L^u(\Omega, \mathcal F,\mathbb P)$, where $u$ is given by $\frac{1}u=\frac{1}{q_0}+\frac{2}{p}+\frac{2}{q}$.

Furthermore, if $2a_0>3-2/s$ then
\begin{equation}\label{13}
\|D^4y\|_{\infty}\leq R_\om n^{-\kappa}
\end{equation}
where $\kappa=2a_0-3-3/s-\delta$ and  $R_\om\geq1$ is 
a random variable
such that $\omega\mapsto R_\om\in L^{u_1}(\Omega,\mathcal F,\mathbb P)$, with $u_1$  defined by $\frac1{u_1}=\frac{4}{q_0}+\frac{3}{s}$. 
\end{lemma}

\begin{proof}[Proof of Lemma \ref{Lemma 4}]
  For every such a  branch $y$ there are inverse  branches $y_i$ of $T_{\sigma^{-i}\omega,\varepsilon}$  such that 
  $$
y=y_n\circ y_{n-1}\circ\ldots\circ y_1.
  $$
  Thus, 
  $$
Dy=F_n\cdot F_{n-1}\cdots F_1 
  $$
  where
   $F_j=D(y_j)\circ y_{j-1}\circ\ldots \circ y_1$. Therefore, 
\begin{equation}\label{D2F}
  D^2y=\sum_{k=1}^n F_{n}\cdot F_{n-1}\cdots F_{k+1}\cdot D(F_k)\cdot F_{k-1}\cdots F_1.
\end{equation}
  Now, using that $\|F_j\|_\infty\leq \gamma_{\sigma^{-j}\omega}^{-1}$ and that 
  $$
\|D(F_k)\|_\infty\leq\|D^2(y_k)\|_\infty\prod_{s=1}^{k-1}\|Dy_s\|_\infty\leq
c(\sigma^{-k}\omega)
  $$
  we see that 
  $$
\|D^2y\|_\infty\leq\prod_{j=1}^n\gamma_{\sigma^{-j}\omega}^{-1}\sum_{k=1}^n \gamma_{\sigma^{-k}\omega}c(\sigma^{-k}\omega)=\prod_{j=1}^n\gamma_{\sigma^{-j}\omega}^{-1}\sum_{k=1}^n\alpha(\sigma^{-k}\omega),
  $$
where $\alpha(\omega)=c(\omega)\gamma_\omega\in L^s(\Omega,\mathcal F,\mathbb P)$ and $s>0$ 
is given by $1/s=1/p+1/q$. By Lemma~\ref{B lemma}, for every $\delta>0$ 
there is a random variable $R\in L^s(\Omega,\mathcal F,\mathbb P)$ such that for $\mathbb P$ a.e. $\om \in \Omega$ and all $n\in \N$ we have $\alpha(\sigma^{-n}\omega)\leq R(\omega)n^{1/s+\delta}$. Therefore, there is a constant $C=C_{s,\delta}>0$ such that $\mathbb P$-a.s. for all $n\geq1$ we have
  \begin{equation}\label{above}
\sum_{k=1}^n\alpha(\sigma^{-k}\omega)\leq CR(\omega)n^{1+1/s+\delta}.      
  \end{equation}
Now \eqref{11} follows from \eqref{gam dec}.



Next, we establish \eqref{12}.
Using~\eqref{D2F} we have that 
  \begin{equation}\label{D 3}
  \begin{split}
D^3y &=\sum_{k=1}^nF_n\cdot F_{n-1}\cdots F_{k+1}\cdot D^2(F_k)\cdot F_{k-1}\cdots F_1    \\  
&\phantom{=}+2\sum_{1\leq i<j\leq n}F_n\cdots F_{j+1}\cdot D(F_j) \cdot F_{j-1}\cdots F_{i+1}\cdot D(F_i)\cdot F_{i-1}\cdots F_1  \\
&=:I_1+I_2.
\end{split}
\end{equation}
  Let us first bound $\|I_1\|_\infty$. 
 We have 
 \begin{equation}\label{F 1}
D(F_k)=\left(D^2(y_k)\circ y_{k-1}\circ\ldots\circ y_1\right)\cdot F_{k-1}\cdots F_1.
 \end{equation}
  Thus, with $G_k :=D^2(y_k)\circ y_{k-1}\circ\ldots\circ y_1$ we have
 \begin{equation}\label{F 2}
 \begin{split}
D^2(F_k) &=\left(D^3(y_k)\circ y_{k-1}\circ\ldots \circ y_1\right)\cdot (F_{k-1}\cdots F_1)^2 \\
&\phantom{=}+\sum_{j=1}^{k-1}G_k\cdot F_{k-1}\cdots F_{j+1}\cdot D(F_j)\cdot F_{j-1}\cdots F_1 \\
&=:J_{1,k}+J_{2,k}.
\end{split}
\end{equation}
Now, using the above upper bounds on $\|F_{j}\|_\infty$ and $\|D(F_j)\|_\infty$ we see that
\begin{equation}\label{J k 1}
\|J_{1,k}\|_\infty\leq c(\sigma^{-k}\om)\prod_{j=1}^{k-1}\gamma_{\sigma^{-j}\om}^{-2}\leq c(\sigma^{-k}\om)   
\end{equation}
and 
\begin{equation}\label{J k 2}
\begin{split}
\|J_{2,k}\|_\infty &\leq c(\sigma^{-k}\om)  \sum_{j=1}^{k-1}c(\sigma^{-j}\om)\prod_{s=1}^{j-1}\gamma_{\sigma^{-s}\om}^{-1}\prod_{v=j+1}^{k-1}\gamma_{\sigma^{-v}\om}^{-1} \\
&\leq c(\sigma^{-k}\om)\prod_{s=1}^{k-1}\gamma_{\sigma^{-s}\om}^{-1}
\sum_{j=1}^{k-1}\alpha(\sigma^{-j}\omega) \\
&\leq Cc(\sigma^{-k}\om)R(\om)k^{1+1/s+\delta},
\end{split}
\end{equation}
where in the last step we used \eqref{above} with  $n=k$.
Putting together the above estimates and using that $\|F_j\|_\infty\leq \gamma_{\sigma^{-j}\om}^{-1}$
 we get that 
\[
\begin{split}
\|I_1\|_\infty &\leq (1+CR(\om))\sum_{k=1}^n k^{1+1/s+\delta}c(\sigma^{-k}\om)\prod_{j=1}^{k-1}\gamma_{\sigma^{-j}\om}^{-1}\prod_{j=k+1}^{n}\gamma_{\sigma^{-j}\om}^{-1} \\
&=(1+CR(\om))\prod_{j=1}^{n}\gamma_{\sigma^{-j}\om}^{-1}\sum_{k=1}^nk^{1+1/s+\delta}\alpha(\sigma^{-k}\om) \\
&\leq  
C R(\om)(1+CR(\om))n^{2+2/s+2\delta}\prod_{j=1}^n\gamma_{\sigma^{-j}\om}^{-1} \\
&\leq C(1+CR(\om))R(\om)E_\om n^{-(a_0-2-2/s-2\delta)},
 \end{split}
\]
 where  the last inequality  uses \eqref{gam dec}. 
 

In order to bound $\|I_2\|_\infty$, using that $\|F_j\|_\infty\leq \gamma_{\sigma^{-j}\om}^{-1}$ and $\|D(F_j)\|\leq c(\sigma^{-j}\om)$, we see that
  $$
\|I_2\|_\infty\leq 2\sum_{1\leq i<j\leq n}\gamma_{\sigma^{-1}\om}^{-1}\cdots \gamma_{\sigma^{-(i-1)}\om}^{-1}  c(\sigma^{-i}\om)\gamma_{\sigma^{-(i+1)}\om}^{-1}\cdots \gamma_{\sigma^{-(j-1)}\om}^{-1}c(\sigma^{-j}\om)\gamma_{\sigma^{-(j+1)}\om}^{-1}\cdots \gamma_{\sigma^{-n}\omega}^{-1}
$$
$$
\leq 2\prod_{k=1}^n\gamma_{\sigma^{-k}\om}^{-1} \sum_{1\leq i<j\leq n}\alpha(\sigma^{-i}\om)\alpha(\sigma^{-j}\om)\leq \prod_{k=1}^n\gamma_{\sigma^{-k}\om}^{-1}\left(\sum_{j=1}^n \alpha(\sigma^{-j}\om)\right)^2 
  $$
  $$
\leq E_\om n^{-a_0}C^2(R(\om))^2n^{2+2/s+2\delta},
  $$
  where in the last inequality we used~\eqref{gam dec} and~\eqref{above}. 
Now \eqref{12} follows from the above estimates on $\|I_1\|_\infty$ and $\|I_2\|_\infty$.

Now we bound $D^4y$. Differentiating both sides of \eqref{D 3} and bounding all the terms by their supremum norm we see that 
$$
\|D^4y\|_\infty\leq 8(L_1+L_2+L_3)
$$
where with $\mathcal I_n=\{1,2,...,n\}$
$$
L_1:=\sum_{j=1}^n\left(\prod_{j<s\leq n}\|F_j\|_\infty\right)\|D^3(F_j)\|_\infty \left (\prod_{1\leq s<j}\|F_j\|_\infty \right),
$$
$$
L_2=\sum_{1\leq i,j\leq n, i\not=j}\left(\prod_{s\in \mathcal I_{n}, s\not=i,j}\|F_s\|_\infty\right)\|D^2(F_j)\|_\infty\|D(F_i)\|_\infty 
$$
$$
L_3=\sum_{1\leq i<j<k\leq n}\left(\prod_{s\in \mathcal I_{n}, s\not=i,j,k}\|F_s\|_\infty\right)\|D(F_i)\|_\infty \|D(F_j)\|_\infty \|D(F_k)\|_\infty. 
$$
Next we estimate $\|D^3(F_j)\|_\infty$. We will use the following abbreviation $F_{a,b}:=F_{b}\cdots F_{a+1}\cdot F_a$.
By differentiating both sides of \eqref{F 2} and bounding each term by its supremum norm
we see that 
\[
\begin{split}
\|D^3(F_k)\|_\infty &\leq \|D^4(y_k)\|_\infty \|F_{1,k-1}\|_\infty^3+ 2\|D^3(y_k)\|_\infty\|F_{1,k-1}\|_\infty\sum_{j=1}^{k-1}\|F_{j+1,k-1}D(F_j)F_{1,j-1}\|_\infty \\
&\phantom{\leq}+\sum_{j=1}^{k-1}\|D(G_k)\|_\infty\|F_{j+1,k-1}\|_\infty\|D(F_j)\|_\infty\|F_{1,j-1}\|_\infty \\
&\phantom{\leq}+
\sum_{i,j\in \mathcal I_{k-1}, i\not=j}\|G_k\|_\infty\|D(F_i)\|_\infty\|D(F_j)\|_\infty\prod_{s\in \mathcal I_{k-1}, s\not=i,j}\|F_s\|_\infty \\
&\phantom{\leq}+\sum_{j=1}^{k-1}\|G_k\|_\infty\|D^2(F_j)\|_\infty\prod_{s\in \mathcal I_{k-1}, s\not=j}\|F_s\|_\infty \\
&=:U_1(k)+U_2(k)+U_3(k)+U_4(k)+U_5(k).
\end{split}
\]
Next, denote $\beta_{a,b}(\om)=\prod_{j=a}^{b}\gamma_{\sigma^{-j}\om}^{-1}$. Recall also the notation $\alpha(\om)=\gamma_\om c(\om)$. 
Then, using that $\|F_j\|_\infty\leq \gamma_{\sigma^{-j}\om}^{-1}$ and \eqref{Pre Adler} we see that
\begin{equation}\label{U1}
U_1(k)\leq c(\sigma^{-k}\om)\left(\beta_{1,k-1}(\om)\right)^3.  
\end{equation}
Moreover, using also that $\|D(F_j)\|_\infty \le c(\sigma^{-j}\om)\beta_{1, j-1}(\om)$ (see~\eqref{F 1}) we see that 
\begin{equation}\label{U2}
U_2(k)\leq 2c(\sigma^{-k}\om)\beta_{1,k-1}(\om) \sum_{j=1}^{k-1}\beta_{j+1,k-1}(\om)(\beta_{1,j-1}(\om))^2 c(\sigma^{-j}\om)
\end{equation}
$$
=2c(\sigma^{-k}\om)(\beta_{1,k-1}(\om))^2 \sum_{j=1}^{k-1}\beta_{1,j-1}(\om)\alpha(\sigma^{-j}\om).
$$
Additionally, using that 
$$
D(G_k)= (D^3(y_k) \circ y_{k-1}\circ \ldots \circ y_1)F_{1,k-1}
$$
we have 
\begin{equation}\label{U3}
U_3(k)\leq c(\sigma^{-k}\om)\beta_{1,k-1}(\om)\sum_{j=1}^{k-1}\beta_{j+1,k-1}(\om)(\beta_{1,j-1}(\om))^2 c(\sigma^{-j}\om)
\end{equation}
$$
=c(\sigma^{-k}\om)(\beta_{1,k-1}(\om))^2\sum_{j=1}^{k-1}\beta_{1,j-1}(\om)\alpha(\sigma^{-j}\om).
$$
Furthermore, 
%using again  \eqref{F 1} to bound $\|D(F_j)\|_\infty$ by $c(\sigma^{-j}\om)\beta_{1,j-1}$ 
we get that
\begin{equation}\label{U4}
\begin{split}
U_4(k) &\leq c(\sigma^{-k}\om)\sum_{1\leq i,j\leq k-1}c(\sigma^{-i}\om)\beta_{1,i-1}(\om)c(\sigma^{-j}\om)\beta_{1,j-1}(\om)\prod_{s\not=i,j, 1\leq s\leq k-1}\gamma_{\sigma^{-s}\om}^{-1} \\
&=c(\sigma^{-k}\om)\beta_{1,k-1}(\om)\sum_{1\leq i,j\leq k-1}\alpha(\sigma^{-i}\om)\beta_{1,i-1}(\om)\alpha(\sigma^{-j}\om)\beta_{1,j-1}(\om) \\
&\leq c(\sigma^{-k}\om)\beta_{1,k-1}(\om)\left(\sum_{i=1}^{k-1}\alpha(\sigma^{-i}\om)\beta_{1,i-1}(\om)\right)^2.
\end{split}
\end{equation}
Finally, using \eqref{J k 1} and \eqref{J k 2} to estimate $\|D^2(F_k)\|_\infty$ we see that
\begin{equation}\label{U5}
U_5(k)\leq C'(1+R(\om))c(\sigma^{-k}\om)\sum_{j=1}^{k-1}\beta_{1,j-1}(\om)\beta_{j+1,k-1}(\om)c(\sigma^{-j}\om)j^{1+1/s+\delta}
\end{equation}
$$
 \leq C'(1+R(\om))c(\sigma^{-k}\om)\beta_{1,k-1}(\om)\sum_{j=1}^{k-1}\alpha(\sigma^{-j}\om)j^{1+1/s+\delta},
$$
where $R(\om)\in L^s$ and $\delta>0$ can be taken to be arbitrarily small.
Using \eqref{gam dec}, the above estimates and that $\alpha(\sigma^{-n}\om)\leq R(\om)n^{1/s+\delta}$, we conclude that  there is a constant $C''>0$ such that
$$
\|D^3(F_k)\|_\infty\leq C''c(\sigma^{-k}\om)\Big(E_\om^3 k^{-3a_0}+E_\om^3 R(\om) k^{-(3a_0-1-1/s-\delta)}+E_\om^3 (R(\om))^2 k^{-(3a_0-2/s-2-2\delta)}
$$
$$
+(1+R(\om))R(\om)E_\om k^{-(a_0-2-2/s-2\delta)} \Big)\leq c(\sigma^{-k}\om)V(\om)k^{-\theta},
$$
where $\theta :=a_0-2-2/s-2\delta$
and $V(\om)\in L^d$, $1/d=3/q_0+2/s$. 

Using the above estimates we see that 
\begin{equation*}
  L_1\leq \beta_{1,n}(\om)V(\om)\sum_{j=1}^n j^{-\theta}\alpha(\sigma^{-j}\om)\leq CE_\om V(\om)R(\om)n^{-(a_0+\theta-1-1/s-\delta)}.  
\end{equation*}
Here we take $\delta$ small enough to ensure that $\theta-1/s-\delta\not=-1$ so that 
$
\sum_{j=1}^n j^{-(\theta-1/s-\delta)}=O(n^{-(\theta-1/s-\delta-1)}).
$



Next, we estimate $L_2$. Using \eqref{J k 1} and \eqref{J k 2} and that $\|D(F_k)\|_\infty\leq c(\sigma^{-k}\om)\beta_{1,k-1}(\om)$, 
we see that 
$$
L_2\leq C(1+R(\om))\sum_{1\leq i,j\leq n}c(\sigma^{-j}\om)j^{1+1/s+\delta}c(\sigma^{-i}\om)\beta_{1,i-1} (\om)\prod_{s\in \mathcal I_n, s\not=i,j}\gamma_{\sigma^{-s}\om}^{-1}
$$
$$
=C(1+R(\om))\beta_{1,n}(\om)\sum_{1\leq i,j\leq n}\alpha(\sigma^{-j}\om)j^{1+1/s+\delta}\alpha(\sigma^{-i}\om)\beta_{1,i-1}(\om)
$$
$$
=C(1+R(\om))\beta_{1,n}(\om)\sum_{i=1}^n\alpha(\sigma^{-i}\om)\beta_{1,i-1}(\om)\sum_{j=1}^n \alpha(\sigma^{-j}\om)j^{1+1/s+\delta}.
$$
Using that $\alpha(\sigma^{-k}\om)\leq R(\om)k^{1/s+\delta}$ and \eqref{gam dec} we conclude that 
$$
L_2\leq C'(1+R(\om))E_\om^2 (R(\om))^2 n^{-(2a_0-\frac 3 s -3-3\delta)}.
$$

To complete the proof we need to estimate $L_3$. Using the upper bounds $\|D(F_k)\|_\infty\leq c(\sigma^{-k}\om)\beta_{1,k-1}(\om)$ and $\|F_k\|_\infty\leq \gamma_{\sigma^{-k}\om}^{-1}$ we see that 
\[
\begin{split}
L_3 &\leq \beta_{1,n}(\om)\sum_{1\leq i,j,k\leq n}\alpha(\sigma^{-i}\om)\alpha(\sigma^{-j}\om)\alpha(\sigma^{-k}\om)\beta_{1,i-1}(\om)\beta_{1,j-1}(\om)\beta_{1,k-1}(\om) \\
&=\beta_{1,n}(\om)\left(\sum_{j=1}^n \alpha(\sigma^{-j}\om)\beta_{1,j}(\om)\right)^3.
\end{split}
\]
Using that $\alpha(\sigma^{-j}\om)\leq R(\om)j^{1/s+\delta}$ and \eqref{gam dec} we conclude that 
$$
L_3\leq CE_\om^4(R(\om))^3 n^{-(4a_0-3-3/s-3\delta)}.
$$
Combining the estimates of $L_1,L_2,L_3$ the proof of \eqref{13} is complete.
\end{proof}

\end{document}
